% !TEX root = ../Projektdokumentation.tex
\section{Fazit}
\label{sec:Fazit}
Mit dem \Fachbegriff{AntispamAdminCenter} wurde eine betriebsnahe
Verwaltungsanwendung entwickelt, die den zuvor manuellen Pflegeprozess für
Cloud-basierte Kunden zentralisiert. Mail- und IP-Freigaben müssen dadurch
nicht mehr auf Shell-Ebene gepflegt werden. Der Nutzen liegt nicht nur in der
Zeitersparnis. Die Daten werden einheitlich gepflegt, Rechte sind klarer
vergeben und Änderungen lassen sich später nachvollziehen.

Das Projekt hat gezeigt, dass auch ein kleinerer Anwendungsfall sorgfältig
umgesetzt werden muss, wenn mehrere Firmen, Rollen und
sicherheitsrelevante Vorgaben beteiligt sind. Wenn die Anwendung später
produktiv genutzt wird, sollte sie an eine eigene Datenbank angebunden werden.
Als nächster Schritt wäre außerdem zu prüfen, ob die Anmeldung an die
bestehende Benutzerverwaltung angebunden werden kann zB. Azure AD .

Aus Sicht des Betriebs ist das Ergebnis eine gute Grundlage, um den Prozess
weiter zu vereinheitlichen. Die wichtigsten Probleme aus der alten Arbeitsweise
wurden reduziert: fehlende Transparenz, uneinheitliche Pflege und die
Abhängigkeit von Shell-Kenntnissen.
